% ---
% id: receiver/ass:nosel
% kind: ass
% title: "Device fixed independently of the adversary"
% summary: >-
%   The receiver configuration and the probe declaration are fixed before the
%   run and cannot be chosen or influenced by the adversary. With [[A-TRUST]]
%   this makes one declared POVM, rather than a family, the right object.
% prominence: standard
% status: accepted
% uses: []
% mentions: [ass:trust]
% assumes: []
% refs: []
% provenance: newly-derived
% short_id: A-NOSEL
% verification: never
% category: receiver
% affects: [secrecy]
% refines: A-TRUST
% user_statement: >-
%   Your receiver's configuration (its components and their settings, which
%   output goes to which detector, and the probe settings) is fixed before the
%   run. An attacker cannot choose or change it, including through properties
%   the model does not describe, such as the wavelength of injected light
%   shifting a beam splitter's ratio. If an attacker can steer the receiver
%   this way, the declared measurement is the wrong one and the result does
%   not apply. This premise is about the configuration of the receiver
%   itself; what may leave or enter the station by other routes is the
%   laboratory premise of this station's party, [[A-LAB]].
% lean:
% ---

Together with Assumption~\ref{ass:trust}, this assumption is what makes a
single POVM, rather than a family, the right object.

\begin{assumption}
\label{ass:nosel}
(i) The component list, its parameters, the port map and the
\tref{probe declaration} are fixed before the protocol run;
(ii) none of them is chosen by the adversary, and none can be influenced by
her, in particular not through a degree of freedom the model does not carry,
such as wavelength.
\end{assumption}
